\section{Problemas resueltos de estimación por intervalos y tamaño muestral}

\subsection{Enunciados de los problemas}

\subsubsection*{1. Nivel Fundamental (Sencillos / Directos)}

\begin{problema}
	\label{prob:n_media}
	Un equipo de investigación desea estimar el tiempo promedio en que un usuario permanece interactuando en una plataforma educativa virtual antes de cerrar sesión. Por datos preliminares en plataformas análogas, se sabe que la desviación estándar de la duración de las sesiones ronda los $\sigma = 18$ minutos.
	\begin{enumerate}
		\item ¿Qué tamaño de muestra se requiere para estimar la media poblacional $\mu$ con un margen de error máximo tolerable de $E = 2$ minutos al nivel de confianza del $95\%$?
		\item Si el presupuesto solo permite monitorear un máximo de $n = 150$ sesiones, ¿cuál sería el margen de error real alcanzado con el $95\%$ de confianza?
	\end{enumerate}
\end{problema}

\begin{problema}
	\label{prob:ic_varianza}
	Un equipo de analítica financiera monitorea la volatilidad diaria de un índice bursátil de tecnología. Durante un mes bursátil de $n = 21$ días hábiles, los rendimientos logarítmicos muestran una desviación estándar muestral de $S = 1.45\%$. Asumiendo que los rendimientos diarios siguen una distribución normal, construya un intervalo de confianza del $90\%$ para la verdadera varianza poblacional $\sigma^2$ y para la desviación estándar $\sigma$.
\end{problema}

\begin{problema}
	\label{prob:dif_medias_z_avanzado}
	En una auditoría de rendimiento en centros de datos, se comparan los tiempos de respuesta del procesamiento de consultas en dos clústeres de servidores independientes. De una muestra aleatoria de $n_1 = 50$ consultas en el Clúster 1 se obtuvo un promedio de $\bar{X}_1 = 120.0$ ms con una desviación estándar muestral $S_1 = 15.0$ ms. En el Clúster 2, una muestra de $n_2 = 60$ consultas arrojó $\bar{X}_2 = 112.0$ ms con $S_2 = 12.0$ ms.
	
	Dado que los tamaños muestrales son grandes ($n_1, n_2 \ge 30$), construya un intervalo de confianza paramétrico del $95\%$ para la diferencia en los tiempos promedio de respuesta $\mu_1 - \mu_2$ usando la distribución normal estándar $Z$.
\end{problema}

\subsubsection*{2. Nivel Operativo (Intermedios / De Desarrollo)}

\begin{problema}
	\label{prob:dif_medias_indep_iguales}
	Un ingeniero de ciencia de datos evalúa el tiempo de latencia (en milisegundos) de dos arquitecturas de bases de datos en la nube (Arquitectura A y Arquitectura B) bajo cargas de trabajo similares. Extrae dos muestras aleatorias independientes con los siguientes resultados:
	\begin{itemize}
		\item \textbf{Arquitectura A:} $n_1 = 15$, media muestral $\bar{X}_1 = 48.5$ ms, desviación estándar $S_1 = 6.2$ ms.
		\item \textbf{Arquitectura B:} $n_2 = 12$, media muestral $\bar{X}_2 = 53.8$ ms, desviación estándar $S_2 = 5.8$ ms.
	\end{itemize}
	Asumiendo que las poblaciones son normales y poseen varianzas iguales ($\sigma_1^2 = \sigma_2^2$), construya un intervalo de confianza del $95\%$ para la verdadera diferencia de latencias medias $\mu_1 - \mu_2$.
\end{problema}

\begin{problema}
	\label{prob:dif_proporciones_ab}
	En una prueba A/B para un sitio web de comercio electrónico, se muestran al azar dos versiones de un formulario de pago a los visitantes. De $n_1 = 800$ usuarios expuestos a la Versión A (diseño clásico), $X_1 = 144$ completaron la compra. De $n_2 = 850$ usuarios expuestos a la Versión B (diseño minimalista de un solo paso), $X_2 = 187$ completaron la compra.
	\begin{enumerate}
		\item Construya un intervalo de confianza del $95\%$ para la verdadera diferencia entre las tasas de conversión poblacionales $p_1 - p_2$.
		\item ¿Justifican los datos la decisión de adoptar definitivamente la Versión B?
	\end{enumerate}
\end{problema}

\begin{problema}
	\label{prob:n_proporcion}
	Una firma analítica proyecta realizar una encuesta nacional para determinar la proporción $p$ de hogares urbanos que utilizan servicios de inteligencia artificial generativa en su vida cotidiana.
	\begin{enumerate}
		\item Calcule el tamaño de muestra mínimo requerido para que la estimación tenga un margen de error no mayor a $3\%$ ($E = 0.03$) con un nivel de confianza del $99\%$, asumiendo que no se dispone de ninguna encuesta previa (escenario más conservador).
		\item Suponga que un estudio piloto anterior sugirió que dicha proporción se sitúa en torno al $20\%$ ($\hat{p}^* = 0.20$). ¿Cuántas encuestas se requerirían bajo este nuevo escenario con el mismo nivel de precisión y confianza? ¿Qué porcentaje de ahorro en el tamaño muestral representa contar con esta información previa?
	\end{enumerate}
\end{problema}

\subsubsection*{3. Nivel Analítico (Avanzados / De Conexión)}

\begin{problema}
	\label{prob:dif_medias_welch}
	En un estudio sobre el rendimiento de dos algoritmos de optimización (Algoritmo 1 y Algoritmo 2) aplicados a redes neuronales convolucionales, se midió el tiempo de convergencia en segundos. Como la variabilidad de uno de ellos parece mucho mayor, no se asumen varianzas iguales ($\sigma_1^2 \neq \sigma_2^2$). Los datos obtenidos son:
	\begin{itemize}
		\item \textbf{Algoritmo 1:} $n_1 = 10$, $\bar{X}_1 = 120.4$ s, $S_1 = 18.5$ s.
		\item \textbf{Algoritmo 2:} $n_2 = 14$, $\bar{X}_2 = 105.2$ s, $S_2 = 7.3$ s.
	\end{itemize}
	Construya un intervalo de confianza del $99\%$ para $\mu_1 - \mu_2$ utilizando la aproximación de Welch-Satterthwaite.
\end{problema}

\begin{problema}
	\label{prob:ic_razon_varianzas}
	Para verificar si dos sensores IoT de temperatura industrial fabricados por distintos proveedores tienen exactamente la misma varianza y precisión en sus mediciones, se toman dos muestras independientes:
	\begin{itemize}
		\item \textbf{Proveedor 1:} $n_1 = 16$ mediciones, con varianza muestral $S_1^2 = 0.36$ $^\circ$C$^2$.
		\item \textbf{Proveedor 2:} $n_2 = 13$ mediciones, con varianza muestral $S_2^2 = 0.12$ $^\circ$C$^2$.
	\end{itemize}
	Construya un intervalo de confianza del $98\%$ para la razón de varianzas poblacionales $\dfrac{\sigma_1^2}{\sigma_2^2}$. ¿Existe evidencia estadística para afirmar que un proveedor tiene mayor variabilidad que el otro?
\end{problema}

\subsubsection*{4. Nivel Desafiante (Expertos / Deducción Profunda)}

\begin{problema}
	\label{prob:ic-desaf-fisher}
	\textbf{Deducción del intervalo asintótico exacto para la correlación de Pearson mediante la transformación de Fisher:}
	El coeficiente de correlación lineal muestral $r$ obtenido en una muestra bivariada normal $(X_i, Y_i)$ de tamaño $n$ posee una distribución en el intervalo $[-1, 1]$ que se vuelve extremadamente asimétrica cuando el verdadero parámetro $\rho \to \pm 1$.
	\begin{enumerate}
		\item Explique teóricamente por qué la técnica del error estándar simétrico $r \pm Z_{\alpha/2}\text{SE}(r)$ produce absurdos matemáticos o cobertura deficiente cuando $|\rho|$ es grande en muestras pequeñas o medianas.
		\item La transformación hiperbólica estabilizadora de varianza de Fisher se define por:
		\begin{align}
			Z(r) = \text{artanh}(r) = \frac{1}{2} \ln\left( \frac{1 + r}{1 - r} \right).
		\end{align}
		Demuestre algebraicamente cómo invertir el intervalo simétrico de confianza asintótico para $Z(\rho)$ en la escala transformada ($Z(r) \pm Z_{\alpha/2}\frac{1}{\sqrt{n-3}}$) usando la función tangente hiperbólica $\tanh(z) = \frac{e^{2z}-1}{e^{2z}+1}$ para obtener los límites explícitos $[L_\rho, \, U_\rho]$ en la escala original del coeficiente de correlación.
		\item En un análisis bioinformático sobre expresión génica colineal en $n = 28$ pacientes, se observó un coeficiente de correlación muestral $r = 0.82$. Calcule exactamente el intervalo de confianza del $95\%$ para $\rho$.
	\end{enumerate}
\end{problema}

\begin{problema}
	\label{prob:ic-desaf-delta-method}
	\textbf{El Método Delta para la inferencia por intervalos sobre cocientes de medias:}
	En economía cuantitativa y ciencia de datos, con frecuencia se desea estimar y construir intervalos de confianza para razones o cocientes de parámetros poblacionales $\theta = \dfrac{\mu_Y}{\mu_X}$ (por ejemplo, elasticidad promedio, eficiencia del gasto o retorno sobre costo), a partir de pares de observaciones independientes $(X_1, Y_1), \dots, (X_n, Y_n)$ donde $\mu_X \neq 0$. Se define el estimador cociente natural:
	\begin{align}
		\hat{\theta} = \frac{\bar{Y}}{\bar{X}}.
	\end{align}
	\begin{enumerate}
		\item Utilizando el desarrollo en serie de Taylor de primer orden (o Método Delta multivariado) alrededor del punto exacto $(\mu_X, \mu_Y)$, demuestre rigurosamente que la varianza asintótica de $\hat{\theta}$ está dada por la expresión:
		\begin{align}
			\Var(\hat{\theta}) \approx \frac{1}{n \mu_X^2} \left[ \sigma_Y^2 - 2\theta \sigma_{XY} + \theta^2 \sigma_X^2 \right],
		\end{align}
		donde $\sigma_{XY} = \cov(X, Y)$ es la covarianza poblacional.
		\item A partir del resultado deducido, exprese de forma sistemática cómo construir el intervalo de confianza asintótico del $(1-\alpha)$ para $\theta = \dfrac{\mu_Y}{\mu_X}$ reemplazando los parámetros desconocidos por los momentos muestrales consistentes ($\bar{X}, \bar{Y}, s_X^2, s_Y^2, s_{XY}$).
	\end{enumerate}
\end{problema}


\subsection{Soluciones detalladas}

\subsubsection*{1. Nivel Fundamental (Sencillos / Directos)}

\begin{solproblema}
	\textbf{Parte 1: Tamaño de muestra requerido ($n$).}
	Para un nivel del $95\%$ de confianza, el valor normal tipificado es $Z_{0.025} = 1.96$. La dispersión conocida es $\sigma = 18$ min y el margen de error tolerable es $E = 2$ min.
	
	Sustituimos directamente en la fórmula elemental del tamaño muestral:
	\begin{align}
		n = \left( \frac{Z_{\alpha/2} \cdot \sigma}{E} \right)^2 = \left( \frac{1.96 \times 18}{2} \right)^2 = (1.96 \times 9)^2 = (17.64)^2 = 311.1696.
	\end{align}
	Dado que los tamaños muestrales corresponden a conteos discretos de sesiones enteras y el margen de error debe estrictamente no exceder los 2 minutos, aplicamos la regla del techo redondeando siempre hacia arriba:
	\begin{align}
		n \geq \lceil 311.17 \rceil = 312 \text{ sesiones.}
	\end{align}
	
	\textbf{Parte 2: Margen de error con presupuesto fijo para $n = 150$.}
	Despejamos $E$ de la fórmula del semiancho del intervalo de confianza para la media:
	\begin{align}
		E = Z_{\alpha/2} \frac{\sigma}{\sqrt{n}} = 1.96 \times \frac{18}{\sqrt{150}} = \frac{35.28}{12.2474} \approx 2.8806 \text{ minutos.}
	\end{align}
	Por consiguiente, al limitar la recolección a $150$ sesiones, la precisión del estudio disminuye arrojando un margen de error incrementado de $2.88$ minutos al $95\%$ de confianza.
\end{solproblema}

\begin{solproblema}
	\textbf{Paso 1: Grados de libertad y cuantiles chi-cuadrada.}
	Los grados de libertad del estadístico varianza muestral son $\nu = n - 1 = 21 - 1 = 20$.
	Para un nivel del $90\%$ de confianza ($\alpha = 0.10$), las colas superior e inferior excluyen un área de $\alpha/2 = 0.05$ y $1-\alpha/2 = 0.95$. Consultando las tablas de la distribución $\chi^2$ con $\nu = 20$:
	\begin{align}
		\chi^2_{0.05, \, 20} &= 31.410, \\
		\chi^2_{0.95, \, 20} &= 10.851.
	\end{align}
	
	\textbf{Paso 2: Intervalo para la varianza $\sigma^2$.}
	El numerador de la suma de cuadrados $(n-1)S^2$ es:
	\begin{align}
		(20)(1.45)^2 = 20 \times 2.1025 = 42.05 \text{ (\%}^2\text{).}
	\end{align}
	Sustituyendo en los límites exactos de la distribución $\chi^2$:
	\begin{align}
		\frac{42.05}{31.410} \le \sigma^2 \le \frac{42.05}{10.851} \implies [1.3387, \, 3.8752] \text{ (\%}^2\text{).}
	\end{align}
	
	\textbf{Paso 3: Intervalo para la desviación estándar $\sigma$.}
	Al ser la raíz cuadrada una función estrictamente creciente para reales positivos, tomamos la raíz en todos los términos:
	\begin{align}
		\sqrt{1.3387} \le \sigma \le \sqrt{3.8752} \implies [1.157, \, 1.969] \%.
	\end{align}
	
	\textbf{Interpretación:} Con un $90\%$ de confianza, la verdadera volatilidad diaria del índice bursátil se encuentra acotada entre $1.16\%$ y $1.97\%$.
\end{solproblema}

\begin{solproblema}
	La diferencia entre los promedios muestrales de latencia es:
	\begin{align}
		\hat{D} = \bar{X}_1 - \bar{X}_2 = 120.0 - 112.0 = 8.0 \text{ ms.}
	\end{align}
	Al ser ambas muestras de gran tamaño ($n_1 = 50, n_2 = 60$), el error estándar de la diferencia se calcula y estandariza sin requerir la suposición de varianzas iguales mediante el Teorema del Límite Central:
	\begin{align}
		\text{SE}(\bar{X}_1 - \bar{X}_2) = \sqrt{\frac{S_1^2}{n_1} + \frac{S_2^2}{n_2}} = \sqrt{\frac{15.0^2}{50} + \frac{12.0^2}{60}} = \sqrt{\frac{225.0}{50} + \frac{144.0}{60}} = \sqrt{4.5 + 2.4} = \sqrt{6.9} \approx 2.6268 \text{ ms.}
	\end{align}
	Para un nivel del $95\%$ de confianza ($Z_{0.025} = 1.96$), el margen de error es:
	\begin{align}
		E = 1.96 \times 2.6268 \approx 5.1485 \text{ ms.}
	\end{align}
	El intervalo de confianza paramétrico al 95\% resulta en:
	\begin{align}
		8.0 - 5.1485 \le \mu_1 - \mu_2 \le 8.0 + 5.1485 \implies [2.85, \, 13.15] \text{ ms.}
	\end{align}
	\textbf{Interpretación:} Con un 95\% de confianza, las consultas en el Clúster 1 toman en promedio entre $2.85$ y $13.15$ milisegundos más que en el Clúster 2, probando de forma estadísticamente robusta que el Clúster 2 es más rápido.
\end{solproblema}

\subsubsection*{2. Nivel Operativo (Intermedios / De Desarrollo)}

\begin{solproblema}
	\textbf{Paso 1: Cálculo de la varianza muestral combinada ($S_p^2$).}
	Al asumir normalidad y homoscedasticidad ($\sigma_1^2 = \sigma_2^2$), obtenemos la varianza agrupada ponderada:
	\begin{align}
		S_p^2 &= \frac{(n_1 - 1)S_1^2 + (n_2 - 1)S_2^2}{n_1 + n_2 - 2} = \frac{(15 - 1)(6.2)^2 + (12 - 1)(5.8)^2}{15 + 12 - 2} \\
		&= \frac{14(38.44) + 11(33.64)}{25} = \frac{538.16 + 370.04}{25} = \frac{908.2}{25} = 36.328 \text{ ms}^2.
	\end{align}
	Por ende, la desviación estándar agrupada es $S_p = \sqrt{36.328} \approx 6.027$ ms.
	
	\textbf{Paso 2: Cuantil crítico en la distribución $t$ con $\nu = 25$.}
	Los grados de libertad totales son $\nu = 15 + 12 - 2 = 25$. Para el 95\% de confianza ($\alpha = 0.05$), el valor en tablas es:
	\begin{align}
		t_{0.025, \, 25} = 2.060.
	\end{align}
	
	\textbf{Paso 3: Construcción del intervalo.}
	La diferencia de medias es $\bar{X}_1 - \bar{X}_2 = 48.5 - 53.8 = -5.3$ ms. El margen de error $E$ se calcula como:
	\begin{align}
		E = t_{0.025, \, 25} \cdot S_p \sqrt{\frac{1}{n_1} + \frac{1}{n_2}} = 2.060(6.027) \sqrt{\frac{1}{15} + \frac{1}{12}} = (12.4156) \sqrt{0.1500} = (12.4156)(0.3873) \approx 4.8086 \text{ ms.}
	\end{align}
	El intervalo del $95\%$ para la diferencia de latencias es por tanto:
	\begin{align}
		-5.3 - 4.809 \le \mu_1 - \mu_2 \le -5.3 + 4.809 \implies [-10.11, \, -0.49] \text{ ms.}
	\end{align}
	\textbf{Interpretación:} Al estar todo el intervalo contenido en la región negativa ($[-10.11, -0.49]$) y no tocar el cero, se confirma al 95\% de confianza que la Arquitectura A tiene una latencia media significativamente inferior (es más veloz) que la Arquitectura B.
\end{solproblema}

\begin{solproblema}
	\textbf{Parte 1: Proporciones y error estándar.}
	Las proporciones empíricas de compra (conversión) son:
	\begin{align}
		\hat{p}_1 &= \frac{144}{800} = 0.18 \quad (18\%), \\
		\hat{p}_2 &= \frac{187}{850} = 0.22 \quad (22\%).
	\end{align}
	La diferencia puntual de tasas es $\hat{p}_1 - \hat{p}_2 = 0.18 - 0.22 = -0.04$. Verificamos que todas las cuentas observadas y esperadas de éxitos y fracasos ($144, 656, 187, 663$) son superiores a $10$, por lo que el uso de la aproximación normal asintótica es totalmente válido.
	
	El error estándar de la diferencia se calcula como:
	\begin{align}
		\text{SE}(\hat{p}_1 - \hat{p}_2) &= \sqrt{\frac{\hat{p}_1(1-\hat{p}_1)}{n_1} + \frac{\hat{p}_2(1-\hat{p}_2)}{n_2}} = \sqrt{\frac{0.18(0.82)}{800} + \frac{0.22(0.78)}{850}} \\
		&= \sqrt{\frac{0.1476}{800} + \frac{0.1716}{850}} = \sqrt{0.0001845 + 0.0002019} = \sqrt{0.0003864} \approx 0.01966.
	\end{align}
	Para el $95\%$ de confianza ($Z_{0.025} = 1.96$), el margen de error $E$ es:
	\begin{align}
		E = 1.96 \times 0.01966 \approx 0.0385 \text{ (o 3.85\%).}
	\end{align}
	El intervalo de confianza del $95\%$ para $p_1 - p_2$ resulta en:
	\begin{align}
		-0.04 - 0.0385 \le p_1 - p_2 \le -0.04 + 0.0385 \implies [-0.0785, \, -0.0015].
	\end{align}
	
	\textbf{Parte 2: Decisión sobre el diseño A/B.}
	El intervalo para $p_1 - p_2$ abarca desde $-7.85\%$ hasta $-0.15\%$. Dado que el límite superior es estrictamente inferior a cero, existe certeza estadística al 95\% de confianza de que la tasa de conversión de la Versión B ($p_2$) supera en forma real a la de la Versión A ($p_1$). Por lo tanto, los datos justifican la adopción inmediata de la Versión B.
\end{solproblema}

\begin{solproblema}
	\textbf{Parte 1: Escenario conservador ($\hat{p}^* = 0.5$).}
	Para el 99\% de confianza, el cuantil normal es $Z_{0.005} = 2.576$. El margen de error admitido es $E = 0.03$. Al carecer de encuestas previas, maximizamos la varianza con $p^*(1-p^*) = 0.5(0.5) = 0.25$:
	\begin{align}
		n_{\text{max}} = \frac{Z_{\alpha/2}^2}{4 E^2} = \frac{(2.576)^2}{4 (0.03)^2} = \frac{6.635776}{4 (0.0009)} = \frac{6.635776}{0.0036} = 1,843.27.
	\end{align}
	Aplicando redondeo superior por conteo entero: $n = \lceil 1,843.27 \rceil = 1,844$ hogares en total.
	
	\textbf{Parte 2: Escenario con información previa ($\hat{p}^* = 0.20$).}
	Aprovechando el estudio piloto, la varianza estimada se reduce a $p^*(1-p^*) = 0.20(0.80) = 0.16$:
	\begin{align}
		n = \left( \frac{Z_{\alpha/2}}{E} \right)^2 p^*(1-p^*) = \left( \frac{2.576}{0.03} \right)^2 (0.16) = (85.8667)^2 (0.16) = (7,373.084)(0.16) = 1,179.69.
	\end{align}
	Redondeando al entero superior: $n = \lceil 1,179.69 \rceil = 1,180$ hogares en la encuesta.
	
	\textbf{Evaluación del porcentaje de ahorro:}
	El ahorro absoluto en encuestas al aplicar la información de ciencia de datos a priori es $1,844 - 1,180 = 664$ hogares. La reducción porcentual de costos de recolección es:
	\begin{align}
		\text{Ahorro} = \frac{664}{1,844} \times 100\% \approx 36.01\%.
	\end{align}
\end{solproblema}

\subsubsection*{3. Nivel Analítico (Avanzados / De Conexión)}

\begin{solproblema}
	\textbf{Paso 1: Grados de libertad de Welch-Satterthwaite ($\nu$).}
	Calculamos las varianzas muestrales divididas entre sus respectivos tamaños de grupo:
	\begin{align}
		\frac{S_1^2}{n_1} = \frac{18.5^2}{10} = \frac{342.25}{10} = 34.225, \quad \frac{S_2^2}{n_2} = \frac{7.3^2}{14} = \frac{53.29}{14} \approx 3.8064.
	\end{align}
	Sustituyendo en la aproximación analítica para varianzas heterogéneas:
	\begin{align}
		\nu &= \frac{\left( \dfrac{S_1^2}{n_1} + \dfrac{S_2^2}{n_2} \right)^2}{\dfrac{(S_1^2 / n_1)^2}{n_1 - 1} + \dfrac{(S_2^2 / n_2)^2}{n_2 - 1}} = \frac{(34.225 + 3.8064)^2}{\dfrac{34.225^2}{10 - 1} + \dfrac{3.8064^2}{14 - 1}} = \frac{(38.0314)^2}{\dfrac{1,171.3506}{9} + \dfrac{14.4887}{13}} \\
		&= \frac{1,446.3874}{130.1501 + 1.1145} = \frac{1,446.3874}{131.2646} \approx 11.0188.
	\end{align}
	Tomando la parte entera conservadora para no sobreestimar la confianza, obtenemos $\nu = 11$ grados de libertad.
	
	\textbf{Paso 2: Valor crítico en la tabla de Student y error estándar.}
	Para el 99\% de confianza ($\alpha = 0.01 \implies \alpha/2 = 0.005$) en $\nu = 11$, el valor en tablas es $t_{0.005, \, 11} = 3.106$.
	El error estándar de la diferencia entre promedios es:
	\begin{align}
		\text{SE}(\bar{X}_1 - \bar{X}_2) = \sqrt{\frac{S_1^2}{n_1} + \frac{S_2^2}{n_2}} = \sqrt{34.225 + 3.8064} = \sqrt{38.0314} \approx 6.1670 \text{ s.}
	\end{align}
	
	\textbf{Paso 3: Construcción e interpretación del intervalo al 99\%:}
	La diferencia muestral observada es $\hat{D} = 120.4 - 105.2 = 15.2$ s. El margen de error $E$ resulta en:
	\begin{align}
		E = 3.106 \times 6.1670 \approx 19.155 \text{ s.}
	\end{align}
	El intervalo al 99\% para $\mu_1 - \mu_2$ es:
	\begin{align}
		15.2 - 19.155 \le \mu_1 - \mu_2 \le 15.2 + 19.155 \implies [-3.955, \, 34.355] \text{ s.}
	\end{align}
	Al cruzar por el cero de negativo a positivo ($[-3.96, 34.36]$), no podemos declarar que el Algoritmo 2 sea significativamente más veloz que el Algoritmo 1 al 99\% de confianza.
\end{solproblema}

\begin{solproblema}
	\textbf{Paso 1: Grados de libertad y puntos críticos de Fisher ($F$).}
	Los grados de libertad del numerador y denominador son $\nu_1 = n_1 - 1 = 15$ y $\nu_2 = n_2 - 1 = 12$.
	Para una confianza del $98\%$ ($\alpha = 0.02 \implies \alpha/2 = 0.01$), el cuantil superior en la tabla $F$ es:
	\begin{align}
		F_{0.01, \, 15, \, 12} = 3.96.
	\end{align}
	Por la propiedad de reciprocidad geométrica para intercambiar grados de libertad, el cuantil inferior es:
	\begin{align}
		F_{0.99, \, 15, \, 12} = \frac{1}{F_{0.01, \, 12, \, 15}} = \frac{1}{3.67} \approx 0.2725.
	\end{align}
	
	\textbf{Paso 2: Cálculo del intervalo para la razón poblacional $\dfrac{\sigma_1^2}{\sigma_2^2}$.}
	El cociente empírico entre las varianzas muestrales observadas es:
	\begin{align}
		\frac{S_1^2}{S_2^2} = \frac{0.36}{0.12} = 3.00.
	\end{align}
	Sustituimos en la fórmula general del intervalo para razón de dispersiones:
	\begin{align}
		\frac{S_1^2 / S_2^2}{F_{0.01, \, 15, \, 12}} \le \frac{\sigma_1^2}{\sigma_2^2} \le \frac{S_1^2 / S_2^2}{F_{0.99, \, 15, \, 12}} \implies \frac{3.00}{3.96} \le \frac{\sigma_1^2}{\sigma_2^2} \le \frac{3.00}{0.2725} \implies [0.7576, \, 11.009].
	\end{align}
	\textbf{Interpretación y conclusión docimásica:} El intervalo del 98\% para el cociente de varianzas se extiende de $0.758$ a $11.009$. Puesto que la razón de equivalencia exacta ($1.0$) yace cómodamente en el interior de los límites empíricos, los datos demuestran que no existe evidencia suficiente al $98\%$ para concluir que el sensor del Proveedor 1 tenga una variabilidad superior o inferior que el del Proveedor 2.
\end{solproblema}

\subsubsection*{4. Nivel Desafiante (Expertos / Deducción Profunda)}

\begin{solproblema}
	\begin{enumerate}
		\item \textbf{Limitaciones teóricas de la aproximación simétrica en $r$:}
		El coeficiente de correlación $r$ está estrictamente delimitado en su dominio al espacio compacto $[-1, 1]$. Cuando la verdadera correlación poblacional es fuerte (por ejemplo, $\rho = 0.85$), la distribución en el muestreo de $r$ se comprime violentamente contra la cota superior $+1$, generando una fuerte asimetría negativa con cola izquierda extendida. Si en este escenario utilizáramos la fórmula ingenua $r \pm Z_{\alpha/2}\text{SE}(r)$, el límite superior del intervalo superaría con frecuencia el valor $+1$ (absurdo axiomático en correlaciones), y la tasa de cobertura real caería drásticamente por debajo del nominal $95\%$ debido a que la varianza exacta del estimador $r$ depende pesadamente del propio parámetro $\rho$: $\Var(r) \approx \dfrac{(1-\rho^2)^2}{n}$.
		
		\item \textbf{Inversión de la transformación hiperbólica de Fisher:}
		Sir Ronald Fisher dedujo que al aplicar al estimador $r$ el arcotangente hiperbólico, la variable transformada $Z(r) = \text{artanh}(r) = \dfrac{1}{2}\ln\left(\dfrac{1+r}{1-r}\right)$ sigue una distribución asintóticamente normal cuya media es $Z(\rho)$ y, lo que es fundamental, \textbf{cuya varianza es estabilizada en la constante independiente del parámetro $\dfrac{1}{n-3}$}:
		\begin{align}
			Z(r) \sim N\left( \frac{1}{2}\ln\left(\frac{1+\rho}{1-\rho}\right), \; \frac{1}{n-3} \right).
		\end{align}
		Por lo tanto, el intervalo asintótico simétrico al $(1-\alpha)$ para el parámetro $Z(\rho)$ en la escala transformada se escribe sin error:
		\begin{align}
			\left[ Z_L, \, Z_U \right] = \left[ Z(r) - Z_{\alpha/2} \frac{1}{\sqrt{n-3}}, \; Z(r) + Z_{\alpha/2} \frac{1}{\sqrt{n-3}} \right].
		\end{align}
		Dado que la función tangente hiperbólica $\tanh(z) = \dfrac{e^{2z}-1}{e^{2z}+1}$ es una función estrictamente monótona creciente y continua para todo $z \in \mathbb{R}$, preserva las desigualdades de probabilidad algebraicas al ser aplicada a todos los miembros:
		\begin{align}
			Z_L \le Z(\rho) \le Z_U \iff \tanh(Z_L) \le \tanh\left(\text{artanh}(\rho)\right) \le \tanh(Z_U) \iff \tanh(Z_L) \le \rho \le \tanh(Z_U).
		\end{align}
		Así se deducen los límites exactos en la escala original:
		\begin{align}
			[L_\rho, \, U_\rho] = \left[ \tanh(Z_L), \; \tanh(Z_U) \right].
		\end{align}
		
		\item \textbf{Cálculo numérico para el estudio bioinformático ($r = 0.82, n = 28$):}
		Evaluamos la transformación de Fisher en $r = 0.82$:
		\begin{align}
			Z(0.82) = \frac{1}{2} \ln\left( \frac{1 + 0.82}{1 - 0.82} \right) = \frac{1}{2} \ln\left( \frac{1.82}{0.18} \right) = \frac{1}{2} \ln(10.1111) = \frac{1}{2}(2.31363) \approx 1.1568.
		\end{align}
		El error estándar constante en la escala $Z$ para $n = 28$ es:
		\begin{align}
			\text{SE}_Z = \frac{1}{\sqrt{n - 3}} = \frac{1}{\sqrt{28 - 3}} = \frac{1}{\sqrt{25}} = \frac{1}{5} = 0.2000.
		\end{align}
		Para 95\% de confianza ($Z_{0.025} = 1.96$), construimos los límites intermedios $Z_L$ y $Z_U$:
		\begin{align}
			Z_L &= 1.1568 - 1.96(0.2000) = 1.1568 - 0.3920 = 0.7648, \\
			Z_U &= 1.1568 + 1.96(0.2000) = 1.1568 + 0.3920 = 1.5488.
		\end{align}
		Aplicando la función de inversión monótona $\tanh(z)$ a cada límite para regresar al espacio $[-1, 1]$:
		\begin{align}
			L_\rho &= \tanh(0.7648) = \frac{e^{2(0.7648)} - 1}{e^{2(0.7648)} + 1} = \frac{e^{1.5296} - 1}{e^{1.5296} + 1} = \frac{4.6163 - 1}{4.6163 + 1} = \frac{3.6163}{5.6163} \approx 0.6439. \\
			U_\rho &= \tanh(1.5488) = \frac{e^{2(1.5488)} - 1}{e^{2(1.5488)} + 1} = \frac{e^{3.0976} - 1}{e^{3.0976} + 1} = \frac{22.1448 - 1}{22.1448 + 1} = \frac{21.1448}{23.1448} \approx 0.9136.
		\end{align}
		El intervalo asintótico exacto al 95\% de confianza para la verdadera correlación poblacional en expresión génica es:
		\begin{align}
			IC_{95\%}(\rho) = [0.6439, \; 0.9136].
		\end{align}
		Obsérvese cómo el intervalo obtenido es lógicamente asimétrico alrededor del estimador puntual ($0.82 - 0.6439 = 0.1761$ de distancia izquierda frente a $0.9136 - 0.82 = 0.0936$ de distancia derecha), respetando perfectamente la frontera natural de $+1$.
	\end{enumerate}
\end{solproblema}

\begin{solproblema}
	\begin{enumerate}
		\item \textbf{Deducción algebraico-diferencial del Método Delta para cocientes:}
		Sea la función objetivo no lineal de dos variables $g(u, v) = \dfrac{v}{u}$, evaluada sobre el vector aleatorio de medias muestrales $(\bar{X}, \bar{Y})$ cuyo verdadero parámetro de esperanza es $\boldsymbol{\mu} = (\mu_X, \mu_Y)$.
		
		Desarrollamos la serie de Taylor bivariada de primer orden de $g(\bar{X}, \bar{Y})$ alrededor del punto paramétrico $(\mu_X, \mu_Y)$:
		\begin{align}
			g(\bar{X}, \bar{Y}) \approx g(\mu_X, \mu_Y) + \left. \frac{\partial g}{\partial u} \right|_{(\mu_X, \mu_Y)} (\bar{X} - \mu_X) + \left. \frac{\partial g}{\partial v} \right|_{(\mu_X, \mu_Y)} (\bar{Y} - \mu_Y).
		\end{align}
		Calculamos las derivadas parciales exactas evaluadas en el punto central:
		\begin{align}
			\frac{\partial g}{\partial u} &= \frac{\partial}{\partial u}\left( \frac{v}{u} \right) = -\frac{v}{u^2} \implies \left. \frac{\partial g}{\partial u} \right|_{(\mu_X, \mu_Y)} = -\frac{\mu_Y}{\mu_X^2} = -\frac{\theta}{\mu_X}. \\
			\frac{\partial g}{\partial v} &= \frac{\partial}{\partial v}\left( \frac{v}{u} \right) = \frac{1}{u} \implies \left. \frac{\partial g}{\partial v} \right|_{(\mu_X, \mu_Y)} = \frac{1}{\mu_X}.
		\end{align}
		Sustituyendo en la expansión lineal tipificada del estimador cociente $\hat{\theta}$:
		\begin{align}
			\hat{\theta} = \frac{\bar{Y}}{\bar{X}} \approx \theta - \frac{\theta}{\mu_X}(\bar{X} - \mu_X) + \frac{1}{\mu_X}(\bar{Y} - \mu_Y).
		\end{align}
		Para obtener la varianza asintótica de $\hat{\theta}$, aplicamos el operador varianza a la aproximación lineal. Recordando que para cualquier combinación lineal $\Var(aW_1 + bW_2) = a^2\Var(W_1) + b^2\Var(W_2) + 2ab\cov(W_1, W_2)$ y que el término constante $\theta$ tiene varianza nula:
		\begin{align}
			\Var(\hat{\theta}) \approx \left( -\frac{\theta}{\mu_X} \right)^2 \Var(\bar{X}) + \left( \frac{1}{\mu_X} \right)^2 \Var(\bar{Y}) + 2\left( -\frac{\theta}{\mu_X} \right)\left( \frac{1}{\mu_X} \right) \cov(\bar{X}, \bar{Y}).
		\end{align}
		Por propiedades del muestreo independiente y apareado i.i.d. de tamaño $n$, las varianzas y covarianzas de los promedios muestrales son:
		\begin{align}
			\Var(\bar{X}) = \frac{\sigma_X^2}{n}, \quad \Var(\bar{Y}) = \frac{\sigma_Y^2}{n}, \quad \cov(\bar{X}, \bar{Y}) = \frac{\sigma_{XY}}{n}.
		\end{align}
		Sustituyendo estas varianzas muestrales en la ecuación y extrayendo el factor común $1 / (n \mu_X^2)$:
		\begin{align}
			\Var(\hat{\theta}) \approx \frac{\theta^2}{\mu_X^2} \left( \frac{\sigma_X^2}{n} \right) + \frac{1}{\mu_X^2} \left( \frac{\sigma_Y^2}{n} \right) - \frac{2\theta}{\mu_X^2} \left( \frac{\sigma_{XY}}{n} \right) = \frac{1}{n \mu_X^2} \left[ \sigma_Y^2 - 2\theta \sigma_{XY} + \theta^2 \sigma_X^2 \right].
		\end{align}
		Queda demostrado formalmente el resultado fundamental del Método Delta para cocientes algebraicos de promedios.
		
		\item \textbf{Procedimiento sistemático de estimación por intervalos en la práctica:}
		Dado que en la realidad empírica de ciencia de datos los parámetros teóricos $(\mu_X, \sigma_X^2, \sigma_Y^2, \sigma_{XY}, \theta)$ se desconocen, para construir el intervalo de confianza $(1-\alpha)$ se sigue este protocolo secuencial:
		\begin{itemize}
			\item Se estiman puntualmente los promedios y varianzas muestrales insesgadas: $\bar{X}, \bar{Y}, s_X^2, s_Y^2$, y se calcula la covarianza muestral de los pares dados:
			\begin{align}
				s_{XY} = \frac{1}{n-1} \sum_{i=1}^n (X_i - \bar{X})(Y_i - \bar{Y}).
			\end{align}
			\item Se evalúa la razón puntual central $\hat{\theta} = \dfrac{\bar{Y}}{\bar{X}}$.
			\item Se sustituyen las estimaciones en la fórmula deducida para obtener el error estándar consistente ($\text{SE}_{\hat{\theta}}$):
			\begin{align}
				\text{SE}(\hat{\theta}) = \sqrt{ \frac{1}{n \bar{X}^2} \left[ s_Y^2 - 2\hat{\theta} s_{XY} + \hat{\theta}^2 s_X^2 \right] }.
			\end{align}
			\item Finalmente, por normalidad asintótica del estimador linealizado al $n \to \infty$, el intervalo de confianza al $(1-\alpha)$ para la razón poblacional $\theta$ se formula como:
			\begin{align}
				IC_{(1-\alpha)}(\theta) = \hat{\theta} \pm Z_{\alpha/2} \cdot \text{SE}(\hat{\theta}) = \frac{\bar{Y}}{\bar{X}} \pm Z_{\alpha/2} \sqrt{ \frac{1}{n \bar{X}^2} \left[ s_Y^2 - 2\hat{\theta} s_{XY} + \hat{\theta}^2 s_X^2 \right] }.
			\end{align}
		\end{itemize}
	\end{enumerate}
\end{solproblema}
